
\documentclass{article}
\usepackage{colortbl}
\usepackage{makecell}
\usepackage{multirow}
\usepackage{supertabular}

\begin{document}

\newcounter{utterance}

\centering \large Interaction Transcript for game `function\_detective', experiment `numbers\_passive\_labeled\_pairs', episode 21 with mistral{-}large{-}3{-}azure.
\vspace{24pt}

{ \footnotesize  \setcounter{utterance}{1}
\setlength{\tabcolsep}{0pt}
\begin{supertabular}{c@{$\;$}|p{.15\linewidth}@{}p{.15\linewidth}p{.15\linewidth}p{.15\linewidth}p{.15\linewidth}p{.15\linewidth}}
    \# & \multicolumn{2}{c}{Player} && \multicolumn{2}{c}{Game Master} \\
    \hline

    \theutterance \stepcounter{utterance}  
    & & & \multicolumn{4}{p{0.6\linewidth}}{
        \cellcolor[rgb]{0.9,0.9,0.9}{
            \makecell[{{p{\linewidth}}}]{
                \texttt{\tiny{[P1$\langle$GM]}}
                \texttt{I have a hidden Python function for you to figure out.} \\
\\ 
\texttt{Your job is to infer what it does and then submit a correct Python implementation.} \\
\\ 
\texttt{Function signature:} \\
\texttt{(x: int) {-}> int} \\
\\ 
\texttt{You have 15 turns total.} \\
\\ 
\texttt{I reveal candidate input{-}output pairs one at a time and tell you whether each one is correct.} \\
\\ 
\texttt{You must make your best SOLVE attempt no later than the second{-}to{-}last turn.} \\
\texttt{Do not wait until the very end.} \\
\\ 
\texttt{You can use these actions:} \\
\texttt{{-} JSON format described at the end of the prompt} \\
\texttt{{-} SOLVE: ```python ... ```} \\
\\ 
\texttt{Do not add import statements in your guess.} \\
\texttt{Anything you might need is already available, so just write the function body in hypothesis and return the answer.} \\
\\ 
\texttt{Preloaded examples:} \\
\texttt{1. ((2, 2), True)} \\
\\ 
\texttt{Scoring:} \\
\texttt{{-} correctness is required for any non{-}zero score} \\
\texttt{{-} in interactive modes, solving with fewer information requests is better} \\
\\ 
\texttt{Input formatting reminders:} \\
\texttt{{-} strings must be quoted} \\
\texttt{{-} booleans must be True/False} \\
\texttt{{-} use commas between multiple values} \\
\texttt{{-} do not write any extra action tags in the same message} \\
\\ 
\texttt{Format example:} \\
\texttt{You write:} \\
\texttt{\{'current\_mode': <probing/confirming/don't know; choose whichever best describes your current state in the task>, 'input\_rationale': <brief description of why you chose this hypothesis>, 'hypothesis': <either the exact string "unknown" if you have no guess yet OR your current best guess as runnable Python matching the given signature, with a one{-}line docstring stating what it does. Never write placeholder code, if you cannot commit to a guess, use "unknown". Match the given signature.} \\
\texttt{e.g.} \\
\texttt{def f(x):} \\
\texttt{	"""Returns the absolute value of x."""} \\
\texttt{	return x if x >= 0 else {-}x} \\
\\ 
\texttt{Use "unknown" if you have no guess.>, 'confidence\_score': <score between 1{-}100. 100 means highest confidence in your hypothesis>\}} \\
\\ 
\texttt{Output only one valid json per turn.} \\
            }
        }
    }
     \\ \\

    \theutterance \stepcounter{utterance}  
    & \multicolumn{4}{p{0.6\linewidth}}{
        \cellcolor[rgb]{0.561,0.737,0.561}{
            \makecell[{{p{\linewidth}}}]{
                \texttt{\tiny{[P1$\rangle$GM]}}
                \texttt{```json} \\
\texttt{\{} \\
\texttt{    "current\_mode": "probing",} \\
\texttt{    "input\_rationale": "The preloaded example (2, 2) suggests that the function might return the input unchanged. This is a simple starting hypothesis to test.",} \\
\texttt{    "hypothesis": "def f(x):\textbackslash{}n    \textbackslash{}"\textbackslash{}"\textbackslash{}"Returns the input integer unchanged.\textbackslash{}"\textbackslash{}"\textbackslash{}"\textbackslash{}n    return x",} \\
\texttt{    "confidence\_score": 50} \\
\texttt{\}} \\
\texttt{```} \\
            }
        }
    }
    & & \\ \\

    \theutterance \stepcounter{utterance}  
    & & & \multicolumn{4}{p{0.6\linewidth}}{
        \cellcolor[rgb]{0.9,0.9,0.9}{
            \makecell[{{p{\linewidth}}}]{
                \texttt{\tiny{[P1$\langle$GM]}}
                \texttt{OUTPUT: ((1000, 1001), False)} \\
            }
        }
    }
     \\ \\

    \theutterance \stepcounter{utterance}  
    & \multicolumn{4}{p{0.6\linewidth}}{
        \cellcolor[rgb]{0.561,0.737,0.561}{
            \makecell[{{p{\linewidth}}}]{
                \texttt{\tiny{[P1$\rangle$GM]}}
                \texttt{```json} \\
\texttt{\{} \\
\texttt{    "current\_mode": "probing",} \\
\texttt{    "input\_rationale": "The previous hypothesis (identity function) was incorrect since (1000, 1001) returned False. The function does not simply return the input. The next guess is that it might return x + 1, as 1000 + 1 = 1001, which matches the output.",} \\
\texttt{    "hypothesis": "def f(x):\textbackslash{}n    \textbackslash{}"\textbackslash{}"\textbackslash{}"Returns the input integer incremented by 1.\textbackslash{}"\textbackslash{}"\textbackslash{}"\textbackslash{}n    return x + 1",} \\
\texttt{    "confidence\_score": 70} \\
\texttt{\}} \\
\texttt{```} \\
            }
        }
    }
    & & \\ \\

    \theutterance \stepcounter{utterance}  
    & & & \multicolumn{4}{p{0.6\linewidth}}{
        \cellcolor[rgb]{0.9,0.9,0.9}{
            \makecell[{{p{\linewidth}}}]{
                \texttt{\tiny{[P1$\langle$GM]}}
                \texttt{OUTPUT: (({-}99, {-}92), True)} \\
            }
        }
    }
     \\ \\

    \theutterance \stepcounter{utterance}  
    & \multicolumn{4}{p{0.6\linewidth}}{
        \cellcolor[rgb]{0.561,0.737,0.561}{
            \makecell[{{p{\linewidth}}}]{
                \texttt{\tiny{[P1$\rangle$GM]}}
                \texttt{```json} \\
\texttt{\{} \\
\texttt{    "current\_mode": "probing",} \\
\texttt{    "input\_rationale": "The previous hypothesis (x + 1) fails for ({-}99, {-}92) since {-}99 + 1 = {-}98 ≠ {-}92. The function is not a simple increment. The difference between {-}99 and {-}92 is +7, while 2 and 2 (from the first example) suggests no change. This hints at a more complex operation, possibly involving digit manipulation or absolute value adjustments.",} \\
\texttt{    "hypothesis": "def f(x):\textbackslash{}n    \textbackslash{}"\textbackslash{}"\textbackslash{}"Returns the input integer with its last digit set to 2 if it was negative, otherwise returns x unchanged.\textbackslash{}"\textbackslash{}"\textbackslash{}"\textbackslash{}n    if x < 0:\textbackslash{}n        return (x // 10) * 10 + 2\textbackslash{}n    return x",} \\
\texttt{    "confidence\_score": 40} \\
\texttt{\}} \\
\texttt{```} \\
            }
        }
    }
    & & \\ \\

    \theutterance \stepcounter{utterance}  
    & & & \multicolumn{4}{p{0.6\linewidth}}{
        \cellcolor[rgb]{0.9,0.9,0.9}{
            \makecell[{{p{\linewidth}}}]{
                \texttt{\tiny{[P1$\langle$GM]}}
                \texttt{OUTPUT: (({-}1, 7), False)} \\
            }
        }
    }
     \\ \\

    \theutterance \stepcounter{utterance}  
    & \multicolumn{4}{p{0.6\linewidth}}{
        \cellcolor[rgb]{0.561,0.737,0.561}{
            \makecell[{{p{\linewidth}}}]{
                \texttt{\tiny{[P1$\rangle$GM]}}
                \texttt{```json} \\
\texttt{\{} \\
\texttt{    "current\_mode": "probing",} \\
\texttt{    "input\_rationale": "The last hypothesis failed for ({-}1, 7). The function does not simply replace the last digit with 2. The output 7 for input {-}1 suggests a relationship with the absolute value or digit sum. Another possibility is that the function returns the sum of the digits of the absolute value of x, but for {-}99 this would be 18, not {-}92. Alternatively, it might return the absolute value plus some transformation. Let's test if it returns the sum of the digits of (x + 10) or similar.",} \\
\texttt{    "hypothesis": "def f(x):\textbackslash{}n    \textbackslash{}"\textbackslash{}"\textbackslash{}"Returns the sum of the digits of (x + 10) if x is negative, otherwise returns x unchanged.\textbackslash{}"\textbackslash{}"\textbackslash{}"\textbackslash{}n    if x < 0:\textbackslash{}n        return sum(int(d) for d in str(abs(x + 10)))\textbackslash{}n    return x",} \\
\texttt{    "confidence\_score": 30} \\
\texttt{\}} \\
\texttt{```} \\
            }
        }
    }
    & & \\ \\

    \theutterance \stepcounter{utterance}  
    & & & \multicolumn{4}{p{0.6\linewidth}}{
        \cellcolor[rgb]{0.9,0.9,0.9}{
            \makecell[{{p{\linewidth}}}]{
                \texttt{\tiny{[P1$\langle$GM]}}
                \texttt{OUTPUT: ((10, 10), True)} \\
            }
        }
    }
     \\ \\

    \theutterance \stepcounter{utterance}  
    & \multicolumn{4}{p{0.6\linewidth}}{
        \cellcolor[rgb]{0.561,0.737,0.561}{
            \makecell[{{p{\linewidth}}}]{
                \texttt{\tiny{[P1$\rangle$GM]}}
                \texttt{```json} \\
\texttt{\{} \\
\texttt{    "current\_mode": "probing",} \\
\texttt{    "input\_rationale": "The function returns the input unchanged for positive numbers (e.g., 2, 10). For negative numbers, it transforms them in a non{-}trivial way ({-}99 → {-}92, {-}1 → 7). The pattern suggests it might be adding the absolute value of the last digit multiplied by some factor. For {-}99: {-}99 + 7 = {-}92; for {-}1: {-}1 + 8 = 7. The added value seems to be (9 {-} last\_digit + 1). Let's test this hypothesis.",} \\
\texttt{    "hypothesis": "def f(x):\textbackslash{}n    \textbackslash{}"\textbackslash{}"\textbackslash{}"If x is negative, adds (9 {-} last\_digit + 1) to x; otherwise returns x unchanged.\textbackslash{}"\textbackslash{}"\textbackslash{}"\textbackslash{}n    if x < 0:\textbackslash{}n        last\_digit = abs(x) \% 10\textbackslash{}n        return x + (9 {-} last\_digit + 1)\textbackslash{}n    return x",} \\
\texttt{    "confidence\_score": 60} \\
\texttt{\}} \\
\texttt{```} \\
            }
        }
    }
    & & \\ \\

    \theutterance \stepcounter{utterance}  
    & & & \multicolumn{4}{p{0.6\linewidth}}{
        \cellcolor[rgb]{0.9,0.9,0.9}{
            \makecell[{{p{\linewidth}}}]{
                \texttt{\tiny{[P1$\langle$GM]}}
                \texttt{OUTPUT: ((99, 100), False)} \\
            }
        }
    }
     \\ \\

    \theutterance \stepcounter{utterance}  
    & \multicolumn{4}{p{0.6\linewidth}}{
        \cellcolor[rgb]{0.561,0.737,0.561}{
            \makecell[{{p{\linewidth}}}]{
                \texttt{\tiny{[P1$\rangle$GM]}}
                \texttt{```json} \\
\texttt{\{} \\
\texttt{    "current\_mode": "probing",} \\
\texttt{    "input\_rationale": "The last hypothesis failed for (99, 100). The function does not simply increment positive numbers. The pattern for negatives seems to involve adding a value based on the last digit, but positives might follow a different rule. For 99, the output is 100 (increment by 1), while for 2 and 10, it's unchanged. This suggests the function might increment numbers ending with 9. Let's test this combined hypothesis.",} \\
\texttt{    "hypothesis": "def f(x):\textbackslash{}n    \textbackslash{}"\textbackslash{}"\textbackslash{}"Increments x by 1 if its last digit is 9 or if x is negative (with last{-}digit adjustment); otherwise returns x unchanged.\textbackslash{}"\textbackslash{}"\textbackslash{}"\textbackslash{}n    if x < 0:\textbackslash{}n        last\_digit = abs(x) \% 10\textbackslash{}n        return x + (9 {-} last\_digit + 1)\textbackslash{}n    elif x \% 10 == 9:\textbackslash{}n        return x + 1\textbackslash{}n    else:\textbackslash{}n        return x",} \\
\texttt{    "confidence\_score": 75} \\
\texttt{\}} \\
\texttt{```} \\
            }
        }
    }
    & & \\ \\

    \theutterance \stepcounter{utterance}  
    & & & \multicolumn{4}{p{0.6\linewidth}}{
        \cellcolor[rgb]{0.9,0.9,0.9}{
            \makecell[{{p{\linewidth}}}]{
                \texttt{\tiny{[P1$\langle$GM]}}
                \texttt{OUTPUT: (({-}1000, {-}993), True)} \\
            }
        }
    }
     \\ \\

    \theutterance \stepcounter{utterance}  
    & \multicolumn{4}{p{0.6\linewidth}}{
        \cellcolor[rgb]{0.561,0.737,0.561}{
            \makecell[{{p{\linewidth}}}]{
                \texttt{\tiny{[P1$\rangle$GM]}}
                \texttt{```json} \\
\texttt{\{} \\
\texttt{    "current\_mode": "confirming",} \\
\texttt{    "input\_rationale": "The latest input ({-}1000, {-}993) fits the pattern: {-}1000 + 7 = {-}993, where 7 is (9 {-} 0 + 1). This confirms the negative case logic. The positive case logic (increment if last digit is 9) also holds for 99 → 100. The remaining uncertainty is whether there are other edge cases, but the pattern seems consistent. I will submit the final solution.",} \\
\texttt{    "hypothesis": "def f(x):\textbackslash{}n    \textbackslash{}"\textbackslash{}"\textbackslash{}"If x is negative, adds (10 {-} last\_digit) to x; if x's last digit is 9, increments x by 1; otherwise returns x unchanged.\textbackslash{}"\textbackslash{}"\textbackslash{}"\textbackslash{}n    if x < 0:\textbackslash{}n        last\_digit = abs(x) \% 10\textbackslash{}n        return x + (10 {-} last\_digit)\textbackslash{}n    elif x \% 10 == 9:\textbackslash{}n        return x + 1\textbackslash{}n    else:\textbackslash{}n        return x",} \\
\texttt{    "confidence\_score": 95} \\
\texttt{\}} \\
\texttt{```} \\
            }
        }
    }
    & & \\ \\

    \theutterance \stepcounter{utterance}  
    & & & \multicolumn{4}{p{0.6\linewidth}}{
        \cellcolor[rgb]{0.9,0.9,0.9}{
            \makecell[{{p{\linewidth}}}]{
                \texttt{\tiny{[P1$\langle$GM]}}
                \texttt{OUTPUT: ((1, 2), False)} \\
            }
        }
    }
     \\ \\

    \theutterance \stepcounter{utterance}  
    & \multicolumn{4}{p{0.6\linewidth}}{
        \cellcolor[rgb]{0.561,0.737,0.561}{
            \makecell[{{p{\linewidth}}}]{
                \texttt{\tiny{[P1$\rangle$GM]}}
                \texttt{```json} \\
\texttt{\{} \\
\texttt{    "current\_mode": "probing",} \\
\texttt{    "input\_rationale": "The last hypothesis failed for (1, 2). The function does not increment positive numbers ending with 9, nor does it leave them unchanged. The output 2 for input 1 suggests a possible pattern where the function returns x + 1 for positive numbers, but this contradicts earlier examples (2 → 2, 10 → 10). Alternatively, it might return the next integer if x is odd. Let's test this hypothesis.",} \\
\texttt{    "hypothesis": "def f(x):\textbackslash{}n    \textbackslash{}"\textbackslash{}"\textbackslash{}"If x is negative, adds (10 {-} last\_digit) to x; if x is positive and odd, returns x + 1; otherwise returns x unchanged.\textbackslash{}"\textbackslash{}"\textbackslash{}"\textbackslash{}n    if x < 0:\textbackslash{}n        last\_digit = abs(x) \% 10\textbackslash{}n        return x + (10 {-} last\_digit)\textbackslash{}n    elif x > 0 and x \% 2 == 1:\textbackslash{}n        return x + 1\textbackslash{}n    else:\textbackslash{}n        return x",} \\
\texttt{    "confidence\_score": 80} \\
\texttt{\}} \\
\texttt{```} \\
\\ 
\texttt{**SOLVE Attempt (Final Submission):**} \\
\texttt{```python} \\
\texttt{def f(x):} \\
\texttt{    """If x is negative, adds (10 {-} last\_digit) to x; if x is positive and odd, returns x + 1; otherwise returns x unchanged} \\
            }
        }
    }
    & & \\ \\

    \theutterance \stepcounter{utterance}  
    & & & \multicolumn{4}{p{0.6\linewidth}}{
        \cellcolor[rgb]{0.9,0.9,0.9}{
            \makecell[{{p{\linewidth}}}]{
                \texttt{\tiny{[P1$\langle$GM]}}
                \texttt{Game Violation: ParseError: Response is not valid JSON.} \\
\texttt{Output ONLY:} \\
\texttt{One JSON with confidence\_score, current\_mode, input\_rationale, and hypothesis keys ...} \\
\texttt{OR} \\
\texttt{SOLVE: ```python ... ```} \\
            }
        }
    }
     \\ \\

    \theutterance \stepcounter{utterance}  
    & \multicolumn{4}{p{0.6\linewidth}}{
        \cellcolor[rgb]{0.561,0.737,0.561}{
            \makecell[{{p{\linewidth}}}]{
                \texttt{\tiny{[P1$\rangle$GM]}}
                \texttt{```json} \\
\texttt{\{} \\
\texttt{    "current\_mode": "confirming",} \\
\texttt{    "input\_rationale": "The function appears to follow these rules: (1) For negative numbers, it adds (10 {-} last\_digit) to x (e.g., {-}99 → {-}99 + 7 = {-}92, {-}1 → {-}1 + 9 = 8, but output was 7, so adjust to (9 {-} last\_digit + 1)). (2) For positive numbers, it increments odd numbers by 1 (e.g., 1 → 2) and leaves even numbers unchanged (e.g., 2 → 2, 10 → 10). This matches all given examples. Final submission.",} \\
\texttt{    "hypothesis": "def f(x):\textbackslash{}n    \textbackslash{}"\textbackslash{}"\textbackslash{}"If x is negative, adds (9 {-} last\_digit + 1) to x; if x is positive and odd, returns x + 1; otherwise returns x unchanged.\textbackslash{}"\textbackslash{}"\textbackslash{}"\textbackslash{}n    if x < 0:\textbackslash{}n        last\_digit = abs(x) \% 10\textbackslash{}n        return x + (9 {-} last\_digit + 1)\textbackslash{}n    elif x > 0 and x \% 2 == 1:\textbackslash{}n        return x + 1\textbackslash{}n    else:\textbackslash{}n        return x",} \\
\texttt{    "confidence\_score": 95} \\
\texttt{\}} \\
\texttt{```} \\
\\ 
\texttt{**SOLVE (Final Answer):**} \\
            }
        }
    }
    & & \\ \\

    \theutterance \stepcounter{utterance}  
    & & & \multicolumn{4}{p{0.6\linewidth}}{
        \cellcolor[rgb]{0.9,0.9,0.9}{
            \makecell[{{p{\linewidth}}}]{
                \texttt{\tiny{[P1$\langle$GM]}}
                \texttt{Game Violation: ParseError: Response is not valid JSON.} \\
\texttt{Output ONLY:} \\
\texttt{One JSON with confidence\_score, current\_mode, input\_rationale, and hypothesis keys ...} \\
\texttt{OR} \\
\texttt{SOLVE: ```python ... ```} \\
            }
        }
    }
     \\ \\

    \theutterance \stepcounter{utterance}  
    & \multicolumn{4}{p{0.6\linewidth}}{
        \cellcolor[rgb]{0.561,0.737,0.561}{
            \makecell[{{p{\linewidth}}}]{
                \texttt{\tiny{[P1$\rangle$GM]}}
                \texttt{```python} \\
\texttt{def f(x):} \\
\texttt{    """If x is negative, adds (10 {-} last\_digit) to x; if x is positive and odd, returns x + 1; otherwise returns x unchanged."""} \\
\texttt{    if x < 0:} \\
\texttt{        last\_digit = abs(x) \% 10} \\
\texttt{        return x + (10 {-} last\_digit)} \\
\texttt{    elif x > 0 and x \% 2 == 1:} \\
\texttt{        return x + 1} \\
\texttt{    else:} \\
\texttt{        return x} \\
\texttt{```} \\
            }
        }
    }
    & & \\ \\

    \theutterance \stepcounter{utterance}  
    & & & \multicolumn{4}{p{0.6\linewidth}}{
        \cellcolor[rgb]{0.9,0.9,0.9}{
            \makecell[{{p{\linewidth}}}]{
                \texttt{\tiny{[P1$\langle$GM]}}
                \texttt{Game Violation: ParseError: Response is not valid JSON.} \\
\texttt{Output ONLY:} \\
\texttt{One JSON with confidence\_score, current\_mode, input\_rationale, and hypothesis keys ...} \\
\texttt{OR} \\
\texttt{SOLVE: ```python ... ```} \\
            }
        }
    }
     \\ \\

    \theutterance \stepcounter{utterance}  
    & \multicolumn{4}{p{0.6\linewidth}}{
        \cellcolor[rgb]{0.561,0.737,0.561}{
            \makecell[{{p{\linewidth}}}]{
                \texttt{\tiny{[P1$\rangle$GM]}}
                \texttt{SOLVE:} \\
\texttt{```python} \\
\texttt{def f(x):} \\
\texttt{    """Returns x with last digit set to 0 if negative, x+1 if positive odd, otherwise returns x unchanged."""} \\
\texttt{    if x < 0:} \\
\texttt{        return (x // 10) * 10} \\
\texttt{    elif x > 0 and x \% 2 == 1:} \\
\texttt{        return x + 1} \\
\texttt{    return x} \\
\texttt{```} \\
            }
        }
    }
    & & \\ \\

    \theutterance \stepcounter{utterance}  
    & & & \multicolumn{2}{p{0.3\linewidth}}{
        \cellcolor[rgb]{0.9,0.9,0.9}{
            \makecell[{{p{\linewidth}}}]{
                \texttt{\tiny{[GM$|$GM]}}
                \texttt{Game Verdict: LOSS} \\
            }
        }
    }
    & & \\ \\

    \theutterance \stepcounter{utterance}  
    & & & \multicolumn{2}{p{0.3\linewidth}}{
        \cellcolor[rgb]{0.9,0.9,0.9}{
            \makecell[{{p{\linewidth}}}]{
                \texttt{\tiny{[GM$|$GM]}}
                \texttt{Incorrect.} \\
\texttt{Failed Cases:} \\
\texttt{Input: {[}{-}99{]} | Expected: {-}92 | Got: {-}100} \\
\texttt{Input: {[}{-}1{]} | Expected: 6 | Got: {-}10} \\
\texttt{Input: {[}99{]} | Expected: 99 | Got: 100} \\
            }
        }
    }
    & & \\ \\

    \theutterance \stepcounter{utterance}  
    & & & \multicolumn{2}{p{0.3\linewidth}}{
        \cellcolor[rgb]{0.9,0.9,0.9}{
            \makecell[{{p{\linewidth}}}]{
                \texttt{\tiny{[GM$|$GM]}}
                \texttt{def add\_seven\_if\_negative(x: int) {-}> int:} \\
\texttt{    if x < 0:} \\
\texttt{        return x + 7} \\
\texttt{    return x} \\
            }
        }
    }
    & & \\ \\

    \theutterance \stepcounter{utterance}  
    & & & \multicolumn{2}{p{0.3\linewidth}}{
        \cellcolor[rgb]{0.9,0.9,0.9}{
            \makecell[{{p{\linewidth}}}]{
                \texttt{\tiny{[GM$|$GM]}}
                \texttt{def f(x):} \\
\texttt{    """Returns x with last digit set to 0 if negative, x+1 if positive odd, otherwise returns x unchanged."""} \\
\texttt{    if x < 0:} \\
\texttt{        return (x // 10) * 10} \\
\texttt{    elif x > 0 and x \% 2 == 1:} \\
\texttt{        return x + 1} \\
\texttt{    return x} \\
            }
        }
    }
    & & \\ \\

    \theutterance \stepcounter{utterance}  
    & & & \multicolumn{2}{p{0.3\linewidth}}{
        \cellcolor[rgb]{0.9,0.9,0.9}{
            \makecell[{{p{\linewidth}}}]{
                \texttt{\tiny{[GM$|$GM]}}
                \texttt{Mode: passive\_labeled\_pairs} \\
\texttt{Accuracy: 0.260} \\
\texttt{Information requests: 7/15} \\
\texttt{Revealed passive examples: 8} \\
\texttt{Efficiency: 0.600} \\
\texttt{Slack: 8} \\
\texttt{Observed{-}pair consistency: 0.500} \\
\texttt{Membership balance: 4 true / 4 false} \\
\texttt{Parse errors: 3} \\
\texttt{Runtime errors: 0} \\
            }
        }
    }
    & & \\ \\

\end{supertabular}
}

\end{document}
